\subsection{Forecast action utility against a trivial always-macro policy}
\label{sec:p1-exposed-action-null}
The previously exposed BNCI2014\_001 nine-person
independent LDA scoring implementation generated
nine \emph{positive} source-only forecasts for the
macro-AUC minus proper-predictive-AUC advantage,
where only seven target advantages were positive.
Thus its $7/9$ sign accuracy equals the trivial
always-positive sign predictor. We additionally
asked whether magnitude thresholds on this \emph{same
frozen prediction vector} would have produced
nontrivial decisions, without any model refit.
Always selecting source macro achieved a descriptive
$+0.01260824$ target mean AUC relative to always
selecting the proper-predictive shrinkage. A
post-outcome subject-specific oracle would have
achieved $+0.01515346$ on these same already
observed targets; its availability is entirely
counterfactual and is not a practical selector.

At an illustrative source-forecast threshold
$\tau=0.005$ AUC, a target-blind policy would
select source macro for five held targets
and proper predictive shrinkage for four.
Its descriptive outcome was $-0.00254522$
AUC relative to \emph{always macro};
two target records were worse than macro
by over $0.005$ AUC. With threshold zero,
all nine forecasts favour macro and the
decision is exactly the same as always macro.
The smaller inspected $\tau=0.0025$ produces
a descriptive gain of only $+0.00015003$
AUC over always macro, but these thresholds
were inspected \emph{after} target outcomes
were known and cannot be called prospective,
optimised or statistically established.
The machine-readable source scalar records
are re-audited with exact decimal arithmetic,
and zero-target feature access does not undo
prior \emph{project-wide outcome exposure}.

This counterexample distinguishes apparently
reasonable forecast error from the decision
value of the prediction. No new shrinkage
algorithm, individual safety policy, or
independent target validation follows.
